\section{28 nm CMOS 外围参考条件与通用开销}

\subsection{估计对象与共同条件的层次}

我们要回答一个参考设计问题：将不同介质的局部阵列配以共同的 28 nm CMOS 外围，在声明逻辑运算、内部展开和更新模式后，能够提供多大的 streaming 输入服务能力与 resident 写入能力。证据可以来自互补的器件、外围及宏级论文，组装后的结果由本文推导。

服务边界包含一个逻辑矩阵、实现其权重所需的局部阵列组合，以及本地输入寄存器、驱动、读出、数字重构、输出寄存器和一个更新控制域。八个权重位平面共同表示一份矩阵；互补存储若必需，也归入这份矩阵的实现。独立服务单元数固定为一，不计跨 bank、plane、die 的复制加速。输入从本地宽向量接口被接收，规定输出在本地结果寄存器就绪时完成；上游传输和下游搬运位于此边界之外。resident 服务则从本地写端口接收事务开始，到相应逻辑状态可计算为止。

参数分为三层。\textbf{A 类共同条件}规定逻辑尺寸、格式、边界、CMOS 数字节拍与局部写入资源；\textbf{B 类路径条件}分别规定 ACIM 的激励与 ADC 资源、DCIM 的数字计算组织；\textbf{C 类介质输入}提供物理响应、专用驱动、写入状态转换、读后恢复及刷新。共同 28 nm 是外围参考选择，原文工艺、电压和操作终点保持原样，器件物理时间也保留其自身条件。

\subsection{v0.1 推荐配置与输出约定}

推荐采用 $n_{\rm out}\times n_{\rm in}=128\times128$ 的逻辑矩阵，streaming 与 resident 均为二进制补码 INT8。沿用 Task 0 的单位，$b_S=b_R=1\,\mathrm{Byte/element}$；位数分别为 $8b_S=8b_R=8$。一个输入向量的 payload 是 128 Byte，整份权重的逻辑容量是 16384 Byte。容量只说明状态规模，一次更新的 $B_R$ 由更新事务实际完成的内容决定。

选择 128 行，是因为库内有直接的 128-cell 负载建模和 128 路输入宏；128 个输出使矩阵与逻辑字节换算简单。\sid{SACIM-03} 的物理 $128\times256$ 二进制阵列只对应 64 个四位权重输出；\sid{SDCIM-01} 的 $128\times128$ bit 阵列对应 16 个八位权重输出。本文的 $128\times128$ 明确指\emph{逻辑元素}，因此二进制映射需 $128\times128\times8=131072$ 个数据 bit，而不是直接照抄原宏的物理尺寸。\cite{CMOS-03,SACIM-03,SDCIM-01}

\begin{table}[htbp]
\centering\small
\begin{tabularx}{\textwidth}{@{}p{27mm} X@{}}
\toprule
项目 & 推荐选择与用途（均为本文参考选择）\\
\midrule
局部逻辑服务 & $128\times128$ INT8 矩阵乘一个 INT8 输入向量；交付 128 个 24-bit 结果寄存器。\\
CMOS 工作条件 & 28 nm，标称参考 0.9 V、25\,$^\circ$C；采用近标称工作模式的工程时隙范围。介质专用电压、电平转换另计。\\
输入与权重展开 & ACIM 每次输入 1 bit，128 路二电平驱动；8 个二进制权重位平面并行，平面内每列累加 128 项。\\
ACIM 读出共享 & 每个位平面 16 个 SAR，共 128 个；每个 ADC 在本平面复用 8 个输出列。每个共享组重新建立、求值和转换。\\
读出与重构 & 名义 10-bit SAR，约 8-bit 有效精度目标，归一到 8-bit 部分和码；16 条输出重构通道，每批用两拍进行位权合并和累加。\\
DCIM 组织 & 输入 1 bit/拍、权重 8 bit 并行；32 个输入项/组，16 个输出通道；输入位、行组、输出组顺序处理。\\
本地写入资源 & 一个更新域，128 路物理编码 bit 数据接口；每事务准备与完成各一拍。脉冲驱动的实际并行能力由介质填写。\\
主时序 & 各必要阶段顺序执行，同一单元不预设求值与写入重叠；向量捕获及最终输出提交合计两拍。\\
\bottomrule
\end{tabularx}
\caption{共同配置提案。ADC 数量与 DCIM 通道数是独立资源选择，未施加等面积约束。结构依据见正文，参数入口见共享 JSON。}
\label{tab:configuration}
\end{table}

这里将逻辑位宽与读出精度分开。对一对二进制输入、权重位平面，128 项部分和为 $0\ldots128$，需要 129 个可区分值，8-bit 码可容纳它们。完整 INT8 运算按位权重构，补码符号位使用负权 $-2^7$。24-bit 容器能容纳 128 项 INT8 乘积的整数和，并为中间重构保留空间。\textbf{DCIM 采用精确整数运算；ACIM 采用经校准的近似部分和及数字重构}。有效精度约 8 bit 是转换器目标，阵列误差和量化误差仍会传至最终结果；24-bit 输出容器只规定交付格式。本轮估计完成该服务的时间，算法精度验证留待需要精确性比较时补做。

\Needspace{5\baselineskip}
选择二电平输入和二进制权重位平面，是为了使公共主参考不依赖某种介质必须具备多级状态或高线性多位 DAC。八个位平面同时参与计算属于逻辑精度的资源成本。若介质需多级 cell、差分 cell 或更小子阵列，应提交等效映射与额外轮数；若另选近似模拟多级权重模式，应单列其精度及基线例外。

\subsection{原始证据如何形成外围范围}

\textbf{ADC 采用完整可复用时隙。} \sid{CMOS-02} 是直接的 28 nm 差分异步 SAR 实测依据，名义 10 bit，在 0.9 V、室温下以 100 MS/s 工作，测试结构含输入 buffer 与参考电路。对应的采样周期为
\[
T_{\rm sample}=\frac{1}{100\times10^6}=10\ns .
\]
这个 10 ns 是相邻采样的周期，覆盖采样与逐次比较过程，并非一次 SAR 比较器决策时间。原文没有给出适配当前 CIM 列负载后的独立转换定值。其 1 MHz 输入条件下正文给出 SNDR 55.13 dB、ENOB 8.86 bit；Fig.~10 的 SNDR 标作 56.91 dB，与正文不一致。奈奎斯特输入下报告 SNDR 51.54 dB、ENOB 8.27 bit，Table~2 同样列出后者。因此本参考使用\emph{约 8 bit 有效精度}的量级，不把名义 10 bit 当作十位无误精度。PVT 表来自后仿真，单独保留。\cite{CMOS-02}

参考时隙 $T_A$ 取 10--50 ns，推荐 20 ns。10 ns 保留独立 SAR 的采样周期锚点；20 ns 为移植到复用 CIM 读出的工程预算；50 ns 留出列选择、输入共模/参考适配、采样与局部锁存的宽裕预算。这些倍数是\emph{本文选择}，并非原文测得的 CIM 降速系数。用名义十位 SAR 得到约八位有效读出时，不再因为最终保留八位而按 $8/10$ 缩短周期。范围也不从各论文 headline 的最大、最小 rate 拼接而成。专用跨阻放大、积分或更慢的阵列建立属于 C 类输入。

每个 ADC 的时隙对应一个输出样本，一批 128 个 ADC 并行给出 128 个码，实际为 16 个逻辑输出各自的 8 个权重位平面部分和。一个输入 bit 需 8 个共享组，完整向量需 64 批。共享因子显式来自 $128/16=8$，而不是 ADC 面积与 cell-column 面积之比。\sid{CMOS-02} 的 0.026 mm$^2$ 活性区还包含 buffer/reference，直接复制 128 套该面积也不是当前布局评估；本轮选择的是转换资源预算，面积优化另议。

\textbf{输入采用数字二电平驱动。} 原文中的独立多位 DAC、数字字线激励和电荷耦合网络承担不同功能。\sid{SACIM-03} 使用 128 个四位输入 DAC，四位译码器选择 16 档电压，并经历 reset、evaluation、ADC readout；其宏在 100 MHz 下运行，提供约 10 ns 的\emph{整体}时序检查。p.~3 和 Table~II 采用 0.9 V，结论写作 1.2 V，本文保留这一差异，只取结构及 10 ns 周期量级。\sid{SACIM-05} 则将八位输入分成四组两位、经四电平输入并行求值，16 项累加的宏级 $T_{AC}$ 为 3.8 ns（0.9 V）或 3.6 ns（1.0 V）。其电压叠加有初始化、采样、叠加三阶段，读出后还进行数字移位累加。69 个 ADC 来自原宏的同位权合并，不是普遍的共享系数。\cite{SACIM-03,SACIM-05}

这些实测宏说明，在相应局部负载、低位输入和专门集成下，输入形成可以包含在数 ns 到十 ns 级宏服务内。本文选 $T_I=2$--10 ns、推荐 5 ns，覆盖输入码选择、分布式缓冲、公共开关复位及二电平激励就绪。2 ns 是简化为一位输入后的积极预算，10 ns 保留较宽的扇出与控制裕量；两端均为工程选择，\emph{没有从宏周期中减去假定 cell 延迟}。128 路输入在八个位平面间广播，使用各平面的本地缓冲；大阵列 RC、介质积分、特殊驱动电压的建立由后续介质增加。当前资料足以支持这一主输入方式，无需补一颗不匹配负载的独立八位 DAC。

\textbf{数字路径采用节拍加完整周期数。} \sid{SDCIM-03} 在 0.8 V 下给出 3 ns 计算时钟，八位输入分八拍，并有 32 项 post-sum。\sid{SDCIM-01} 报告 1.1 V 下约 360 MHz（Fig.~10 标 363 MHz），即约 2.8 ns/拍；其 $128\times16$ INT8 VMM 实际需 64 拍，包含输入位展开和空间分组。\sid{SDCIM-05} 的 INT4 一 MAC 拍、INT8 两 MAC 拍；Table~III 的 0.9 V accurate INT4/INT8 access time 为 1.0/2.1 ns，作为更并行结构的交叉依据。该文 FP8 的正文与表格互换了 2.5/2.7 ns 的 AP/AC 标签，故本文不据此选择 INT8 时钟。\cite{SDCIM-01,SDCIM-03,SDCIM-05}

共同数字节拍取 $T_D=2$--10 ns，推荐 5 ns（200 MHz）。2 ns 是具备相应并行与流水资源时的积极设计端；5 ns 围绕几个 ns 的硅测时序留出移植余量；10 ns 为较大局部归约、布线和控制负载保留空间。文献中的更低电压模式可明显慢于此范围，本参考固定近标称模式。来自 0.8、1.0、1.1 V 的锚点不被宣称为固定 0.9 V 的 PVT 保证，也不作电压或工艺线性缩放。

DCIM 一组运算的公共数字阶段包含输入位选择、局部乘法/归约、部分和更新和寄存控制；阵列读侧物理阶段另有定义。ACIM 一批的数字重构分两拍：第一拍合并八个位平面的正负位权，第二拍按输入位权更新 16 个输出累加器。这里的通道数与拍数是资源预算，后续实现若需要更深的归约流水，应显式增加拍数。数字节拍也用于公共命令与完成接口，统一传播其不确定性。

\textbf{写入只共享低压接口及控制预算。} \sid{CMOS-03} 的 28 nm SRAM 分析采用 TT、1 V、50 FO4 约 1 ns 周期、25 FO4 字线脉宽和 15 fF 的 128-cell 位线负载，并以读/写成功及连续操作终点定义周期；这些是模型条件。\sid{CMOS-07} 的 28 nm eFlash 平台 SRAM 有 48-bit 字、1024 字、列 MUX4；Table~II 给出典型 read access 1.54 ns 与 455 MHz 时钟，后者约 2.20 ns。两者共同支持数 ns 级局部普通读写/控制量级，同时显示 access、WL 脉冲与可重复周期是不同量。该 eFlash 工艺具有较高阈值，其数据不被推广为全部逻辑 compiler 的定值。\cite{CMOS-03,CMOS-07}

参考采用 128-bit 物理编码数据接口，结构锚点来自 \sid{SDCIM-01}，而 \sid{SDCIM-05} 的普通端口只有 32 bit；二者都没有独立闭合的写周期。本文选择每事务准备与完成各一拍，$T_W=2T_D=4$--20 ns，推荐 10 ns。第一拍可接收一个 128-bit 数据批，额外数据批各占 $T_D$；page 接口若命令与数据分开，全部数据拍另计。该预算覆盖命令、数据锁存、低压驱动接口和完成状态接收，物理编程脉冲、读回 sensing/ADC、verify 比较和必要高压产生按实际序列加入。若原文给出完整普通写周期并已覆盖这些功能，采用完整周期替代相应块，保持原来已包含的边界。

\Needspace{24\baselineskip}
\subsection{核心范围表与完整服务开销}

% Generated by scripts/check_shared.py --emit.
\begingroup\small
\begin{longtable}{@{}>{\raggedright\arraybackslash}p{23mm}>{\raggedright\arraybackslash}p{49mm}>{\raggedright\arraybackslash}p{24mm}>{\raggedright\arraybackslash}p{58mm}@{}}
\caption{共同外围能力范围。每项按自身粒度使用，完整服务的次数由结构推导。}\label{tab:periphery}\\
\toprule 模块 & 服务粒度 & 范围/推荐 & 依据与用途\\\midrule\endfirsthead
\toprule 模块 & 服务粒度 & 范围/推荐 & 依据与用途\\\midrule\endhead
输入形成与公共复位 $T_I$ & 每次求值的公共复位、128 路二电平输入与本地缓冲。 & 2--10 ns\newline 推荐 5 ns & \sid{SACIM-03}, \sid{SACIM-05}, \sid{CMOS-03}。\newline 更改多位输入或扇出负载时核对驱动与精度；专用建立另计。\\
采样、转换与局部读出 $T_A$ & 一颗 ADC 一个样本的完整可复用时隙；并行颗数属于结构参数。 & 10--50 ns\newline 推荐 20 ns & \sid{CMOS-02}, \sid{SACIM-03}, \sid{SACIM-05}。\newline 名义 10 bit、有效约 8 bit；含 mux/采样/转换/局部锁存，跨批保持另计。\\
数字服务节拍 $T_D$ & 一个局部数字阶段；拍数由运算、归约与通道资源推导。 & 2--10 ns\newline 推荐 5 ns & \sid{SDCIM-01}, \sid{SDCIM-03}, \sid{SDCIM-05}。\newline 参考 200 MHz；共同近标称预算，所选负载下需相应资源。\\
写入准备与完成控制 $T_W$ & 每事务准备和完成两拍；编码数据装入另按接口规则计。 & 4--20 ns\newline 推荐 10 ns & \sid{CMOS-03}, \sid{CMOS-07}, \sid{SDCIM-01}, \sid{SDCIM-05}。\newline 首数据批可与命令同拍；完整更新周期已含接口时替代。\\
\bottomrule\end{longtable}
\endgroup


表中的四块足以组织主参考。ACIM 一次物理步骤经历输入形成、阵列求值、转换、两拍数字重构及必要恢复；$t_{m,A}$ 汇集其中公共预算以外的阵列建立/积分、专用驱动和读后恢复时间，各段按实际操作顺序执行。每个 ADC 共享组重新求值，因而不依赖所有模拟结果长期保持。八个输入 bit 乘八个共享组，共 64 步，向量捕获和最终输出提交另计两拍。因此\emph{纯公共外围占用}为
\begin{equation}
T_{S,A,\rm peri}=64(T_I+T_A+2T_D)+2T_D .
\label{eq:acim_peri}
\end{equation}
固定配置下，短、参考、长时隙情景分别为 1028、2250、5140 ns；展示为约 1--6\,$\mu$s，参考约 2.3\,$\mu$s。它是完整逻辑服务所需的公共外围占用，介质时间仍须加入，未产生任何介质的最终 $\rho$。

DCIM 的 8 个输入 bit、4 个行组、8 个输出组共 256 个计算步骤。每步一拍公共数字阶段，边界两拍，所以纯公共外围占用为 $258T_D=516$--2580 ns，展示约 0.5--3\,$\mu$s，参考约 1.3\,$\mu$s。这些完整服务量级比某些已发表宏的单周期长，主要来自所选逻辑尺寸、位展开和 16 通道复用；原宏的更大并行、模拟位权合并或不同累加维度并未被暗中移植进来。

本配置让计数清楚，但它不是所有介质的唯一合理实现。最值得审阅的是 128 个 ADC 的资源预算、16 条 DCIM 输出通道，以及主参考每个共享组重新求值的组织。这些选择会系统改变服务次数；冻结后，各介质沿用同一基线或明确列出例外。
